\section{Discussion}
\label{sec:discussion}

\paragraph{What the shift means.}
The single most important result is that the 2022 collapse is structural: a classifier
tells 2018 from 2022 Albanian students almost perfectly from background data, and the
variables that moved most are process variables: teacher support fell while
reported home possessions rose. This is not the signature of a household-wealth shock;
it is consistent with a disruption to schooling itself (the pandemic and its long
tail), and it explains why models trained on the pre-crisis trajectory transfer only
moderately to 2022. A ``weaker cohort'' framing would misdirect policy toward
uniform remediation; a structural-shift framing points at the school process.

\paragraph{Why the ceiling matters.}
The instinct on hitting a plateau is to tune harder or add capacity. Here that instinct
is wrong twice over. Tuning did nothing; capacity (a neural net, a kernel machine, a
four-model stack) did nothing. The missing ingredient was a \emph{feature family}, school
socioeconomic composition, and once added, the plateau moves and the model
class that wins changes. But the new plateau at $\sim$0.78 is then real, corroborated
by a multilevel ICC of $\approx$0.31 and by a direct school-questionnaire linkage that
adds nothing. The lesson for educational data mining is methodological: distinguish a
\emph{feature} ceiling (movable, and diagnosable by what a richer feature family buys)
from a \emph{data} ceiling (a property of the measurement instrument), and use
convergent evidence rather than a single learning curve to tell them apart. The
distinction travels: the same feature ceiling appears in all nine systems and the amount
each can move it scales with its between-school inequality ($r\approx0.80$ with the
ICC), and the data ceiling survives the obvious escape route, adding PISA's
within-test process data lifts AUC to $\sim$0.86 but almost entirely through a
cognitive-process signal endogenous to the score, leaving the deployable ceiling at
$\sim$0.78. A plateau, then, can be moved by a new \emph{feature family} or by a new
\emph{modality}, but only by data that a school could actually hold before the test.

\paragraph{Shared drivers, distinctive breadth.}
Albania's risk mechanism is not exotic. Math anxiety and school composition drive risk
across the nine countries; what distinguishes Albania is that low proficiency is
\emph{broad}, sitting on a flat SES gradient and a low floor. This reframes the policy
conversation away from targeting the poor (whom a flat gradient says are not
disproportionately the whole problem) toward system-wide levers.

\paragraph{On causality, and the screener.}
Every model here is associational and fit on a single cross-section. The clearest
illustration is anxiety: lower reported math anxiety co-occurs with higher scores, but
it is partly a \emph{consequence} of doing well, so a counterfactual that ``reduces
anxiety'' in the model is not a defensible intervention estimate. We therefore ship the
risk screener, calibrated probabilities, per-student SHAP drivers, and the fairness
audit, as a \emph{triage and communication} aid, not a lever simulator. The fairness
result reinforces the point: a naive threshold over-flags the poorest quintile, so the
tool is only defensible as a way to widen support; where equal false-positive treatment
across SES is required, group-specific thresholds close the gap from 0.63 to 0.003 at a
measured recall cost, an equity choice we surface rather than bury. The one causal
question we could pose from observational data, whether the 2019 earthquake localised
the collapse, we tested and answered in the negative: a difference-in-differences whose
placebo fails identifies the loss as national, a well-identified null that guards
against reading a convenient single-region story into a country-wide shock.

\paragraph{Limitations.}
The student questionnaire is the only instrument here; we have no teacher or parent
files, and the school questionnaire, while linked, adds no signal beyond composition.
Albania lacks ESCS in 2012 and 2015, limiting its own cross-cycle socioeconomic
comparison (the peer-benchmarked and 2018-vs-2022 analyses are unaffected). Design-based
standard errors are available only for the 2015--2022 cycles; the fixed-width 2009/2012
cycles fall back to an imputation-only SE. The forecast rests on five cycles and one
structural break, which is why we report scenarios rather than a number. Finally, the
$\sim$0.78 ceiling is specific to background-questionnaire prediction of a binary
Level-2 target; we tested one richer modality, PISA's within-test computer-based
process data, and found it moves the plateau only through a signal endogenous to the
score, so genuinely exogenous process data that a school holds before the assessment
(attendance, formative assessment, longitudinal records) remains the untested route and
the natural next step. The difference-in-differences likewise licenses only a negative
causal claim; a design that could license a positive policy lever still requires data we
do not have.
